%% PART II: PROPOSED METHODS AND SETUP
\section*{Proposed Methods}

PrexSyn generates an initial synthesizable candidate for each property specification. Five modification methods then produce structural variants from this seed. All variants, along with a repeated-PrexSyn-sampling baseline, pass through identical synthesizability and property-conservation filters. The baseline tests whether simply drawing more samples from the generative model matches or exceeds the value of a post-generation modification layer.

\textit{CReM} \cite{bib7} modifies molecules by swapping chemical fragments (side chains, ring systems) while respecting chemical validity rules, using a fragment database mined from ChEMBL. A tunable context radius controls how conservatively CReM edits: radius 1 restricts changes to small R-group substitutions (e.g., replacing a methyl group with a chlorine atom); radius 3 permits ring-system and scaffold-level replacements. This mirrors how medicinal chemists optimize leads by iteratively modifying functional groups to improve potency, selectivity, or pharmacokinetics. We benchmark CReM at context radii 1, 2, and 3 to characterize the locality/quality tradeoff within a single method.

\textit{mmpdb} \cite{bib8} mines matched molecular pair transformations: single-site substitution rules extracted from pairs of real compounds that differ by exactly one structural change (e.g., methyl $\rightarrow$ chloro, phenyl $\rightarrow$ pyridyl). Each rule carries empirical property-change annotations, enabling selection of transformations likely to preserve the desired profile. It is the standard tool for conservative, chemist-interpretable local edits in medicinal chemistry.

\textit{JT-VAE} \cite{bib9} encodes molecules into a continuous mathematical space (a learned latent representation) by decomposing each molecule's graph into a tree of chemically meaningful substructures (a junction tree). New variants are generated by perturbing the latent vector and decoding, producing edits spanning medium-to-global structural scales. This enables scaffold hopping, the ability to replace entire core structures while preserving the pharmacological profile, a strategy valued in drug discovery for accessing new chemical series and navigating intellectual property constraints.

\textit{LibINVENT} \cite{bib10,bib11} preserves a molecular scaffold (the core ring system and connectivity) extracted from PrexSyn's output while exploring alternative R-group decorations via reinforcement learning steerable toward target properties.

\textit{ReactEA} \cite{bib13} applies evolutionary search over 22,949 biochemical reaction rules. Starting from a seed molecule, it generates variants by applying reaction transformations iteratively and selecting for desired properties. Every output carries an explicit synthetic route, making it the most synthesis-constrained method in the benchmark.

These five methods span the rule-based/learned axis and the local-to-global edit-scale axis (mmpdb $\rightarrow$ CReM $\rightarrow$ JT-VAE, LibINVENT, ReactEA). Selection was informed by a systematic review of $\sim$50 candidate methods, prioritizing coverage across edit scales and paradigms. No prior work systematically benchmarks modification methods as post-generation complements to synthesis-aware generation. This study provides the first such comparison, including a complementarity analysis identifying which combinations maximize total coverage.

The complementarity analysis also provides the empirical basis for composing effective operators into multi-stage pipelines: chaining local edits, fragment substitutions, and latent-space exploration with synthesis filtering at every stage. No public system yet integrates these operators under unified synthesis constraints; our rankings and complementarity matrices provide the foundation for such integration in both manual medicinal chemistry workflows and automated agent-driven pipelines.
